\documentclass[11pt,a4paper]{article}
\usepackage[margin=22mm,headheight=15pt]{geometry}
\usepackage{fontspec,xeCJK}
\setmainfont{Noto Serif}
\setsansfont{Noto Sans}
\setmonofont{DejaVu Sans Mono}[Scale=0.83]
\setCJKmainfont{Noto Serif CJK SC}
\usepackage{amsmath,amssymb,mathtools,booktabs,longtable,array,xcolor,fancyhdr,hyperref}
\usepackage{xurl}
\setlength{\emergencystretch}{2em}
\definecolor{ink}{HTML}{193248}
\definecolor{accent}{HTML}{156B72}
\hypersetup{colorlinks=true,linkcolor=ink,urlcolor=accent,citecolor=accent,pdftitle={Order-15 third-root maximal determinant: independent audit},pdfauthor={Independent computational audit}}
\pagestyle{fancy}\fancyhf{}\fancyhead[L]{\small\sffamily Order 15 over third roots of unity}\fancyhead[R]{\small\sffamily Independent audit}\fancyfoot[C]{\thepage}
\setlength{\parindent}{0pt}\setlength{\parskip}{6pt}
\newcommand{\muT}{\mu_3}\newcommand{\B}{B}\newcommand{\HH}{H^{*}}\newcommand{\QQ}{\mathbb Q}\newcommand{\ZZ}{\mathbb Z}\newcommand{\tr}{\operatorname{tr}}\newcommand{\diag}{\operatorname{diag}}\newcommand{\spec}{\operatorname{spec}}
\newcommand{\file}[1]{\path{#1}}
\begin{document}
\thispagestyle{empty}
{\sffamily\color{accent}\large COMPUTATIONAL VERIFICATION REPORT}\par\vspace{10mm}
{\Huge\sffamily\color{ink} Order-15 maximal determinant\\[3mm]over third roots of unity}\par\vspace{5mm}
{\Large\sffamily Independent audit and reconstructed certificates}\par
{\large 17 September 2026}\par\vspace{8mm}
\hrule\vspace{5mm}
{\Large\bfseries 中文结论：全局最优性尚未验证}\par
本次复核保留原始压缩包的内容，恢复纪录矩阵，并重新编写、实际运行了一组精确验证程序。已完成全部 $Q\le81$ 的允许壳层，共19个；全部 $Q\ge171$ 由迹—方差界排除，$Q\not\equiv0,6\pmod9$ 由同余排除。部分低能量证明使用已发表的三元等距码定理；该外部分类没有在本次重新运行。

此外，$Q=99$ 的三个局部交织问题已独立恢复全部解并验证矛盾；$Q=105,114,123,132,141,150,159,168$ 的 $(14,1)$ 型内部无孤立顶点分支，全部36种内部能量分配均已取得纪录级上界。这些是完整的局部分支结果，仍不能替代每个壳层的完整证明。

纪录矩阵的三条精确计算路径一致给出
\[
 |\det H_0|^2=\B=277868041444786176=2^{22}3^{20}\cdot19,
 \qquad Q(H_0)=105.
\]
它给出确定的下界。剩余19个允许壳层仍有未完成的证明义务；原 ledger 中的 \texttt{CLOSED} 已与本次证据严格区分。因此，这份文件不得作为“全局最大值已精确确定”的证明提交。

\vfill
{\small\sffamily\color{ink} Delivery scope: preserved inputs, new exact-arithmetic verifiers, execution logs, a 38-shell coverage ledger, and explicit unresolved obligations. No missing repository certificate is mocked.}
\newpage

\section{Statement and scope of the audit}
Let $\omega^2+\omega+1=0$, let $\muT=\{1,\omega,\omega^2\}$, and put
\[
 D_3(15)=\max_{H\in\muT^{15\times15}} |\det H|^2,
 \quad G=HH^*,\quad K=H^*H,\quad
 Q=\sum_{i<j}|G_{ij}|^2.
\]
Both Gram matrices have diagonal 15 and the same energy $Q$. A strict counterexample to the record means $\det G>\B$. This terminology concerns an unproved upper bound; it does not assert that any such matrix exists.

The new computations establish that no strict counterexample has $Q\le81$ or $Q\ge171$. The assertion for $Q\le81$ uses the external theorem $B_3(15,10)=12$ in the cases indicated below. Independently of that external code theorem, the arithmetic/color reductions and the nine shells $Q=0,9,\ldots,72$ are checked by the supplied exact computations. The four shells $6,15,24,33$ are also excluded directly by the color/energy reduction. The remaining six low shells $42,51,60,69,78,81$ use the cited code theorem.

A second result covers every $(14,1)$ color partition with no internal support isolates at $Q=105,114,123,132,141,150,159,168$. Here the upper bound is at most $\B$, with equality possible in the friendship-support case at $Q=105$. A third result exhaustively verifies the specified $A/B$ local obstructions at $Q=99$. Neither result closes all other color and support types at those energies.

The global statement $D_3(15)=\B$ remains \emph{unverified by this delivery}. The unresolved energies are
\[
\begin{split}
&87,90,96,99,105,108,114,117,123,126,\\
&132,135,141,144,150,153,159,162,168.
\end{split}
\]
The package contains the original sixteen archive files, new reconstructions, and machine-readable results. It does not contain a complete checkout of the original working repository. In particular, historical tail closures have not been promoted to verified global statements without replaying their upstream reductions.

\section{Exact arithmetic and the benchmark}
\subsection{Arithmetic representation and cross-checks}
An Eisenstein number $a+b\omega$ is represented by $(a,b)$, with rational coordinates when division is required. Multiplication, conjugation and norm are
\[
(a,b)(c,d)=(ac-bd,\ ad+bc-bd),\qquad
\overline{(a,b)}=(a-b,-b),\qquad N(a,b)=a^2-ab+b^2.
\]
All comparisons deciding a certificate are integer or rational comparisons. Floating-point values in JSON and this report are display-only approximations.

The implementation includes fraction-free Bareiss elimination and a separate rational Gaussian elimination. Sixty deterministically generated matrices of orders 1 through 6 are checked against both methods and against an independent real-integer representation: multiplication by $a+b\omega$ acts through
\[
 \begin{pmatrix}a&-b\\b&a-b\end{pmatrix},
 \qquad \det\mathcal R(A)=N(\det A).
\]
Exact inverse identities and characteristic-polynomial evaluations are also checked. Tree determinant dynamic programming is compared directly with Bareiss elimination for all nonisomorphic trees on 2 through 9 vertices. These tests are in \file{programs/verify_arithmetic_tests.py}; their counts and outcomes are recorded separately from the mathematical shell claims.

\subsection{Recovered witness}
The exponent matrix in \file{data/benchmark.json} was recovered from the repository benchmark data at commit \texttt{007af7d8f7f644396aee27dd87c68c8f4169efd5}. The original data blob is identified in the provenance metadata. The local JSON adds source information and is not claimed to be byte-identical to that blob.

The three exact paths give
\[
 \det H_0=604661760+241864704\omega,\qquad
 \det(H_0H_0^*)=N(\det H_0)=\B.
\]
The row and column energies are both 105. The row support degrees are fourteen occurrences of 2 and one occurrence of 14, giving the friendship support. Consequently
\[
 \sqrt{\B}=2^{11}3^{10}\sqrt{19},\qquad
 \frac{\sqrt{\B}}{15^{15/2}}=0.7965900126038639\ldots.
\]
This is a verified witness and hence a lower bound on $D_3(15)$. The equality with the unknown global optimum is not inferred from it.

\section{Global reductions}
\subsection{Congruence and norm conditions}
For two rows let $n_0+n_1+n_2=15$ count the three exponent differences. Their inner product is
\[
 z=(n_0-n_2)+(n_1-n_2)\omega,
 \qquad \delta=n_1+2n_2\pmod3.
\]
Writing $a=n_0-n_2$ and $b=n_1-n_2$ gives $a+b=15-3n_2$ and $b\equiv\delta\pmod3$. Thus
\[
 N(z)=(a+b)^2-3ab\equiv3\delta^2\pmod9.
\]
The program also checks this identity on all 136 count triples. The row color is the sum of its exponents modulo 3. Same-color edges have squared norm divisible by 9; different-color edges have squared norm 3 modulo 9 and are nonzero. For color-class sizes $a+b+c=15$,
\[
 Q\equiv3(ab+ac+bc)\equiv0\ \text{or}\ 6\pmod9.
\]
There are exactly 38 eligible energies between 0 and 168.

Subtracting the first row from each other row shows that $(1-\omega)^{14}$ divides $\det H$. Therefore $3^{14}$ divides $\det G$. Also, any prime $p\equiv2\pmod3$ occurs to an even exponent in the nonzero norm $\det G=N(\det H)$. This latter necessary condition is used as an exact factorization sieve, never as evidence of realizability.

\subsection{Trace--variance envelope}
For a positive semidefinite $d\times d$ matrix of trace $15d$ and off-diagonal energy $q$, its eigenvalues have mean 15 and squared deviations summing to $2q$. An interior maximum of their product has at most two distinct eigenvalues, by Lagrange multipliers. Boundary products vanish. For multiplicity $m$ the stationary product is
\[
 P_{d,m}(q)=\bigl(15+(d-m)t\bigr)^m\bigl(15-mt\bigr)^{d-m},
 \quad t=\sqrt{\frac{2q}{dm(d-m)}}.
\]
On the positive branch this product decreases with $t$. Replacing $t$ by the certified rational lower bound
\[
 t_- =10^{-15}\left\lfloor10^{15}\sqrt{\frac{2q}{dm(d-m)}}\right\rfloor
\]
therefore gives an upper bound. Taking the maximum over $m$ yields the rational envelope $U_d(q)$ used in the code. Each branch decreases as $q$ increases, so the maximum is nonincreasing as well.

The exact computation verifies
\[
 U_{15}(171)<\B,\qquad U_{15}(171)/\B\approx0.9944244061.
\]
This excludes the entire range $Q\ge171$, not just the single value 171. Exact numerators and denominators are retained in \file{results/foundations.json}.

\subsection{Color partitions}
If a color class of size $a$ is internally orthogonal, the minimum cross energy is $3a(15-a)$. A Schur complement and AM--GM give
\[
 \det G\le15^a\left(\frac{75-a}{5}\right)^{15-a}.
\]
For every $3\le a\le12$, the right-hand side is strictly below $\B$. Hence each such class of a strict counterexample contributes at least one internal norm-9 edge. Combining this with $Q\le168$ leaves ten sorted partitions. Further Schur bounds exclude largest classes of sizes 10, 11 and 12 and leave only
\[
 (15,0,0),\quad(14,1,0),\quad(13,2,0),\quad(13,1,1).
\]
For an active class of size $d$, internal energy $e$, largest-eigenvalue cap $L$, and cross energy at least $c$, the bound is
\[
 U_d(e)\left(15-\frac{c}{L(15-d)}\right)^{15-d}.
\]
The verified $(d,e,L,c)$ values are $(10,9,18,150)$; for $d=11$, $e=9,18,27$ with $L=20,21,23$ and $c=132$; and for $d=12$, $e=9,18,27,36,45,54$ with $L=18,77/4,21,22,241/10,25$ and $c=108$. The low-energy caps follow from the finite norm catalogue and small support graphs; the larger ones follow from $\lambda_{\max}(E)^2\le2e(d-1)/d$. Small unweighted graph caps are checked exactly, including the four-disjoint-edge case beyond the seven-vertex atlas.

\section{Reconstructed low-energy proofs}
\subsection{Isolates and the two-sided projection condition}
Put $E=G-15I$ and $F=K-15I$. The identity $EH=HF$ is central. If $F$ has an isolated support vertex, the corresponding column of $H$ is a full-support vector in $\ker E$. A leaf of the support of $E$ would force one coordinate of that vector to vanish. Therefore an isolate on one side forbids a leaf on the other.

If $F$ has $r$ support isolates, the corresponding columns form a matrix $Y$ with $Y^*Y=15I_r$ and range in $\ker E$. Hence
\[
 \frac1{15}YY^*\preceq P_{\ker E},\qquad
 \min_i(P_{\ker E})_{ii}\ge\frac r{15}.                 \tag{1}
\]
This is a two-sided necessary condition because $E$ and $F$ have the same characteristic polynomial. It is stronger than testing singularity of the active support alone.

\subsection{The nine shells $Q=0,9,\ldots,72$}
For these energies the color reduction forces both Grams to be all-same-color. Up to $Q=63$, at most seven edges cannot cover 15 vertices, so both supports have isolates. At $Q=72$, a support without isolates must be $P_3\sqcup6K_2$ with norm-9 edges. Its Gram determinant is $3105\cdot216^6$, which fails the norm sieve: the primes 5 and 23 have odd valuations. Thus both supports again have isolates, and all active components have minimum degree at least two.

Set $k=Q/9\le8$ and divide $E$ by 3. A connected active component has at most eight edges and eight vertices. The seven-vertex graph atlas plus $C_8$ is exhaustive in this range: an eight-vertex minimum-degree-two graph with at most eight edges is a cycle. Edge norm units are drawn from the exact catalogue; under this budget only 1, 3 and 4 can occur. Sixth-root diagonal switching fixes spanning-tree phases, and every remaining associate class and cycle phase is enumerated. This switching is used only for spectral and projection necessary conditions, which it preserves.

For each component the program computes the characteristic polynomial and the exact kernel projection. Components with a zero projection diagonal are discarded. Component assemblies are filtered by determinant, divisibility and norm. Assemblies with a common full characteristic polynomial are then coupled by (1), iterating until no further state can be removed. When states with identical spectral data are merged, the larger projection minimum is retained; this only relaxes the necessary condition. Every remaining set is empty. No missing historical support catalogue or external equidistant-code theorem is used for these nine shells.

\subsection{The small mixed shells}
An orthogonal set of cube-root rows of length 15 gives a ternary equidistant code of distance 10. The published theorem $B_3(15,10)=12$ \cite{TB} therefore forbids 13 mutually orthogonal rows. This external classification is not rerun here.

A color class of size $a$ with $m$ internal support edges has an independent set of size at least $a-m$. The code theorem forces at least $a-12$ edges when $a>12$. Together with the color congruence and cross energy, this closes $Q=42,51$; $Q=6,15,24,33$ already fail the minimum mixed cross energy of 42. At $Q=60,69,78$, both sides must have color type $(14,1)$.

Let $C$ be the nonisolated part of the internal 14-row support, and let $I$ be its set of $k$ isolates. Here $k\ge6$. The full $15$-eigenspace has zero hub coordinate and coordinate support contained in $I\cup\operatorname{supp}\ker C$. For the opposite Gram, the hub coordinate is likewise absent. Intertwining then forces the cross entries on $I$ into one cube-root orbit: a column of $H$ at that opposite hub is orthogonal to the hyperplane $u_I^\perp$ supported on $I$, and so is proportional to $u_I$. All entries of that column have modulus one and belong to $\muT$.

There are at most two excess cross-energy units, fewer than $k$. The common squared norm of these isolated cross entries must consequently be 3. After cube-root normalization the isolated cross entries are identical. If the opposite $15$-eigenspace is supported on at most $s$ coordinates, at least $15-s$ columns are identical on $I$ up to a scalar cube root. Positivity of the residual Gram matrix gives
\[
 15I_k-(15-s)J_k\succeq0,\qquad k(15-s)\le15.          \tag{2}
\]
The complete finite enumeration contains all internal supports with two to four energy units, every permitted phase, and every active cross allocation. It uses $s\le k+|\operatorname{supp}\ker C|$, a safe upper estimate. The actual enumerated state counts are 36, 1008 and 25416 for $Q=60,69,78$, respectively. After the exact norm sieve and spectral deduplication, (2) removes all states. The output explicitly labels the external code-theorem dependency.

\section{The reconstructed $Q=81$ boundary}
The code theorem removes the mixed $(13,1,1)$ alternative, so both Grams are all-same-color. A no-isolate support has exactly nine norm-9 edges: a heavier edge would reduce the edge count to at most seven, insufficient to cover 15 vertices. Such a graph is either a forest with six components, each of size at most five, or $K_3\sqcup6K_2$. All forest spectra fail the required determinant/norm filters. The surviving no-isolate spectra are the four real-cycle-product possibilities for the triangle, and none has eigenvalue 15.

Therefore, if either support has an isolate, both do. The preceding kernel-projection method extends to nine energy units. For completeness beyond the seven-vertex atlas, an eight-vertex nine-edge minimum-degree-two connected graph has degree excess two. It is a figure-eight, a barbell, or a theta graph. The program explicitly generates all possibilities, together with $C_8$ and $C_9$; it also includes both associate classes of norm-7 edges. The three norm-surviving isolate states are eliminated by (1).

It remains to exclude $K_3\sqcup6K_2$. Every pair block has spectrum $\{18,12\}$. When the triangle cycle product has real part $\pm1/2$ after normalization, its characteristic polynomial is coprime to the pair polynomial. The corresponding $3\times2$ block of $H$ would vanish by intertwining, which is impossible. Positive real cycle product gives a triangle eigenvalue 21, supported on three coordinates on each side. The $3\times3$ block of $H$ would then need squared operator norm at least 21, exceeding its Frobenius norm squared 9.

For negative real cycle product, the triangle spectrum is $\{18,18,9\}$. The localized eigenvalue 9 saturates the same $3\times3$ Frobenius bound. Its block must have rank one. Since its entries are cube roots, row and column cube-root gauges make this block $J_3$ and both triangle Gram blocks $18I_3-3J_3$. Each pair can first be gauged to off-diagonal $\pm3$. A negative pair has a symmetric eigenvector of eigenvalue 12; its two entries against every triangle column would need to sum to zero. Two cube roots cannot do so. All six pair blocks are thus positive.

Each $2\times2$ pair-to-pair block has the form
\[
 \begin{pmatrix}a&b\\b&a\end{pmatrix},\qquad a,b\in\muT.
\]
In symmetric and antisymmetric pair coordinates, write its entries as $S=a+b$ and $D=a-b$. The antisymmetric $6\times6$ block satisfies $DD^*=D^*D=12I_6$. Every row and column has four nonzero entries of squared norm 3 and two zero entries. Correspondingly, an entry of $S$ is $2a$ at a zero of $D$ and is a negative cube root otherwise.

The triangle-to-pair columns and pair-to-triangle rows are balanced cube-root triples. Their Gram conditions force three of each of the two cube-root orbits. Pair gauges and permutations make the two groups identical within each orbit. The cross-Gram equations then say that every row and column of $S$ sums to zero within each group of three. A group cannot contain two zeros of $D$, because $2a+2b-c=0$ has no cube-root solution. Since there are two zeros in total, each group contains exactly one. Also $2a-b-c=0$ forces $a=b=c$.

It follows that each $3\times3$ block of $S$ is $\alpha(3P-J_3)$ for a cube root $\alpha$ and a permutation matrix $P$. Orthogonality between the two row groups would require
\[
 \alpha(9R-3J_3)+\beta(9T-3J_3)=0,
 \qquad \alpha,\beta\in\muT,\quad R,T\in S_3.
\]
All $3^2\cdot6^2=324$ possibilities are checked exactly and none vanishes. Equivalently, on the zero-sum plane a permutation would have to be a scalar equal to the negative of a cube root; the only scalar action of a permutation there is the identity. This completes the no-isolate boundary. The finite checks and the graph/spectrum reductions are in \file{programs/verify_q81.py}.

\section{Independently recovered $Q=99$ local certificates}
The historical final types under examination are
\[
 A=4K_2^-\sqcup K_3^+\sqcup C_4^+,
 \qquad B=3K_2^-\sqcup K_3^+\sqcup T_6,
\]
where $T_6$ has a positive edge between its centers and two negative leaf edges at each center. Both specified full Grams have characteristic polynomial
\[
 (x-21)^2(x-18)^4(x-15)^2(x-12)^6(x-9).
\]
Their $15$-eigenspaces are supported on the four cycle vertices or the four leaves, respectively. This makes the original residual shape conditions necessary: columns outside the relevant local block must be orthogonal to that localized eigenspace.

For each local pair the new verifier constructs the integer linear operator for $LZ-ZR=0$, row-reduces it over $\QQ$, assigns cube roots to the free matrix-entry coordinates, and reconstructs every pivot coordinate. It then checks membership in $\muT$ and compares the entire resulting set with the original parameterization.
\begin{center}
\begin{tabular}{lrrrrr}\toprule
Pair & Variables & Rank & Free & Cube-root solutions & Compatible residuals\\\midrule
$A/B$ &24&18&6&144&0\\
$A/A$ &16&10&6&225&0\\
$B/B$ &36&28&8&729&0\\\bottomrule
\end{tabular}
\end{center}
The raw free-coordinate enumerations have sizes $3^6$, $3^6$ and $3^8$. This validates completeness of the \emph{local} parameterizations and their contradictions. It does not validate an exhaustive reduction of all $Q=99$ matrices to these two global types, nor the separate mixed-color size-13 cases. Those remain explicit obligations.

\section{Record-level component bounds for $(14,1)$}
Assume the 14-row same-color internal support has no isolates. Let its connected components have sizes $v_i\ge2$ and energy units $u_i\ge v_i-1$, so $\sum v_i=14$ and internal energy $e=9\sum u_i$. Write $c=Q-e\ge42$ for the cross energy to the remaining row.

For a general component use determinant upper bound $U_{v_i}(9u_i)$ and largest-eigenvalue cap
\[
 L_i=15+\sqrt{\frac{18u_i(v_i-1)}{v_i}}.
\]
When $u_i=v_i-1$, the component must be a norm-9 tree. All edge phases can be gauged away, and the determinant is exactly that of $15I+3A_T$. Bipartiteness also gives the sharper cap $L_i=15+3\sqrt{u_i}$. Rational upper brackets for every square root preserve the direction of the Schur bound.

The minimum squared norm on every cross entry is 3. Allocating all excess cross energy to the component with the largest cap only weakens the bound, and gives
\[
 \det G\le\prod_i d_i\left(15-\sum_i\frac{3v_i}{L_i}
                  -\frac{c-42}{\max_iL_i}\right),       \tag{3}
\]
with zero as an upper bound when the displayed Schur allowance is nonpositive. Here $d_i$ is the general determinant envelope or the exact maximum over trees of that order.

The tree determinant is computed by its matching polynomial,
\[
 \det(15I+3A_T)=\sum_j(-9)^j m_j(T)15^{|T|-2j},
\]
using exact rooted-tree dynamic programming. All nonisomorphic trees of orders 2 through 14 are enumerated. At order 14, there are 3159 trees; the unique determinant maximizer is $P_{14}$, with value $16802420983158456$. The same enumeration records the unique path maximizer at each smaller order.

All integer component abstractions are enumerated, including unrealizable ones; this enlarges the feasible set safely. Bound (3) is at most $\B$ for every internal energy allocation at each of the eight eligible energies $105,114,\ldots,168$. There are 36 allocations in total. At $Q=105$ the no-isolate support is $7K_2$, and (3) yields
\[
 216^7\left(15-\frac{42}{18}\right)=216^7\frac{38}{3}=\B.
\]
At the other seven energies the computed maxima are strictly below $\B$. This result does not rely on an unverified general ``isolate exclusion'' lemma; its no-isolate hypothesis remains explicit.

\subsection{Why the old $Q=150$ test was insufficient}
The inspected historical script compares its branches to the old $Q=153$ trace envelope. Its six size-14 expressions use
\[
 (e,c,L)=(63,87,25),(72,78,127/5),(81,69,261/10),
\]
\[
 (90,60,27),(99,51,138/5),(108,42,563/20)
\]
and evaluate $U_{14}(e)(15-c/L)$. Five of the six expressions are greater than $\B$, despite all being below the old envelope. Their ratios to $\B$ are approximately
\[
 1.00484,\quad1.00383,\quad1.00457,\quad1.00587,\quad1.00307,\quad0.999981.
\]
This is a certificate-strength gap, not a counterexample. The new component proof repairs this particular no-isolate branch. It still does not supply the missing size-13 state-pairing and isolate arguments, so the whole $Q=150$ shell is not declared closed.

\section{Reproducibility and remaining obligations}
The entry point \file{run_all.py} executes every available verifier in a fresh subprocess and saves stdout, stderr, exit codes and timings. \texttt{--audit-only} requests completion of the audit. Without that option, unresolved global obligations cause exit code 2 even when every available local test passes. Source hashes are checked before execution; Python optimization mode, which disables assertions, is rejected.

The pinned numerical dependencies are SymPy 1.13.1 and NetworkX 3.5. Graph-atlas completeness and nonisomorphic-tree generation are trusted library inputs; this delivery is not a formal proof assistant development and does not independently formalize those graph generators. The external code-classification theorem is identified rather than silently assumed to have been recomputed.

The main missing obligations are the record-level size-13 defect and local inverse arguments; completeness of the large support/phase catalogues; the cubic-majorant, high-triangle and Fourier branches; the $Q=144,150$ state-pairing certificates; and the upstream sparse/spectral coverage of the historical tail closures. The detailed per-shell list is in \file{report/remaining_obligations.md}. The aggregate code never replaces a missing proof by a boolean or a fixed expected answer.

All reported fractions, characteristic polynomials, enumeration counts and partial bounds are saved in \file{results/}. The original archive's ledger remains preserved under \file{original/}, but its labels do not control the new status table. A claim of global maximality requires every unresolved row below to acquire a complete record-level proof.

\appendix
\section{Complete eligible-shell coverage table}
\textbf{Closed} means a complete reconstructed shell argument passed. \textbf{Closed, cited input} additionally uses $B_3(15,10)=12$. \textbf{Unverified} means that a complete global argument for that shell has not been recovered and checked. It does not assert falsity or exhibit a counterexample.

\begin{longtable}{@{}r p{31mm} p{102mm}@{}}
\toprule
$Q$ & Audit status & Evidence / remaining obligation \\
\midrule
\endfirsthead
\toprule
$Q$ & Audit status & Evidence / remaining obligation \\
\midrule
\endhead
\midrule\multicolumn{3}{r}{\small Continued on next page}\\
\endfoot
\bottomrule
\endlastfoot
0 & Closed & Exact no-leaf catalogue, norm sieve and two-sided kernel projection. \\
6 & Closed & Color reduction and minimum cross energy 42. \\
9 & Closed & Exact no-leaf catalogue, norm sieve and two-sided kernel projection. \\
15 & Closed & Color reduction and minimum cross energy 42. \\
18 & Closed & Exact no-leaf catalogue, norm sieve and two-sided kernel projection. \\
24 & Closed & Color reduction and minimum cross energy 42. \\
27 & Closed & Exact no-leaf catalogue, norm sieve and two-sided kernel projection. \\
33 & Closed & Color reduction and minimum cross energy 42. \\
36 & Closed & Exact no-leaf catalogue, norm sieve and two-sided kernel projection. \\
42 & Closed (cited input) & Color/energy bound plus B3(15,10)=12. \\
45 & Closed & Exact no-leaf catalogue, norm sieve and two-sided kernel projection. \\
51 & Closed (cited input) & Color/energy bound plus B3(15,10)=12. \\
54 & Closed & Exact no-leaf catalogue, norm sieve and two-sided kernel projection. \\
60 & Closed (cited input) & Full small mixed support/cross enumeration and eigenprojection PSD; B3(15,10)=12. \\
63 & Closed & Exact no-leaf catalogue, norm sieve and two-sided kernel projection. \\
69 & Closed (cited input) & Full small mixed support/cross enumeration and eigenprojection PSD; B3(15,10)=12. \\
72 & Closed & Exact no-leaf catalogue, norm sieve and two-sided kernel projection. \\
78 & Closed (cited input) & Full small mixed support/cross enumeration and eigenprojection PSD; B3(15,10)=12. \\
81 & Closed (cited input) & Isolate support catalogue and exact K3+6K2 block obstruction; B3(15,10)=12. \\
87 & Unverified & Missing: Record-level size13 defect/local inverse reductions and other remaining support branches; a lower historical comparison target is insufficient. \\
90 & Unverified & Missing: Record-level size13 defect/local inverse reductions and other remaining support branches; a lower historical comparison target is insufficient. \\
96 & Unverified & Missing: Record-level size13 defect/local inverse reductions and other remaining support branches; a lower historical comparison target is insufficient. \\
99 & Unverified & Partial: All AB/AA/BB local mu3 intertwiners independently recovered; 144/225/729 solutions, zero compatible residuals. Missing: Exhaustive reduction to A/B and mixed-color branches, especially size13 internal energy9/18. \\
105 & Unverified & Partial: (14,1) no-internal-isolate branch completely bounded by B over all component allocations.; An explicit matrix attains B; the 7K2 friendship support family has exact upper bound B. Missing: Complete record-level proof for every color, support, and phase case. \\
108 & Unverified & Missing: Record-level size13 defect/local inverse reductions and other remaining support branches; a lower historical comparison target is insufficient. \\
114 & Unverified & Partial: (14,1) no-internal-isolate branch completely bounded by B over all component allocations. Missing: Record-level size13 defect/local inverse reductions and other remaining support branches; a lower historical comparison target is insufficient. \\
117 & Unverified & Missing: Record-level size13 defect/local inverse reductions and other remaining support branches; a lower historical comparison target is insufficient. \\
123 & Unverified & Partial: (14,1) no-internal-isolate branch completely bounded by B over all component allocations. Missing: Record-level size13 defect/local inverse reductions and other remaining support branches; a lower historical comparison target is insufficient. \\
126 & Unverified & Missing: Complete polynomial-majorant/high-triangle orbits, Fourier branches and mixed-color record-level reductions. \\
132 & Unverified & Partial: (14,1) no-internal-isolate branch completely bounded by B over all component allocations. Missing: Record-level size13 defect/local inverse reductions and other remaining support branches; a lower historical comparison target is insufficient. \\
135 & Unverified & Missing: Complete polynomial-majorant/high-triangle orbits, Fourier branches and mixed-color record-level reductions. \\
141 & Unverified & Partial: (14,1) no-internal-isolate branch completely bounded by B over all component allocations. Missing: Record-level size13 defect/local inverse reductions and other remaining support branches; a lower historical comparison target is insufficient. \\
144 & Unverified & Missing: Complete K6 support catalogue and record-level state-pairing certificates. \\
150 & Unverified & Partial: (14,1) no-internal-isolate branch completely bounded by B over all component allocations. Missing: Record-level (13,2) local Schur/state-pairing and (14,1) isolate exclusion; the no-isolate branch has now been reconstructed. \\
153 & Unverified & Missing: Full upstream sparse/spectral/support reductions of the tail certificate, including all-same/mixed residuals. The historical closure claim has not been fully replayed here. \\
159 & Unverified & Partial: (14,1) no-internal-isolate branch completely bounded by B over all component allocations. Missing: Full upstream sparse/spectral/support reductions of the tail certificate, including all-same/mixed residuals. The historical closure claim has not been fully replayed here. \\
162 & Unverified & Missing: Full upstream sparse/spectral/support reductions of the tail certificate, including all-same/mixed residuals. The historical closure claim has not been fully replayed here. \\
168 & Unverified & Partial: (14,1) no-internal-isolate branch completely bounded by B over all component allocations. Missing: Full upstream sparse/spectral/support reductions of the tail certificate, including all-same/mixed residuals. The historical closure claim has not been fully replayed here. \\
\end{longtable}


\section{Execution record}
\small
\begin{longtable}{@{}p{48mm}p{93mm}@{}}
\toprule
Module & Fresh replay result \\
\midrule
\endfirsthead
\toprule
Module & Fresh replay result \\
\midrule
\endhead
\bottomrule
\endfoot
\texttt{verify\_arithmetic\_tests} & Arithmetic cross-checks PASS: 60 matrices; 94 trees \\
\texttt{verify\_foundations} & Foundations PASS: det pair (604661760, 241864704) B 277868041444786176 Q 105 tail/B 0.9944244061347672 \\
\texttt{verify\_low\_energy} & Q 0 VERIFIED\_CLOSED assemblies 1 norm survivors 0 final 0; Q 9 VERIFIED\_CLOSED assemblies 0 norm survivors 0 final 0; Q 18 VERIFIED\_CLOSED assemblies 0 norm survivors 0 final 0; Q 27 VERIFIED\_CLOSED assemblies 0 norm survivors 0 final 0; Q 36 VERIFIED\_CLOSED assemblies 1 norm survivors 0 final 0; Q 45 VERIFIED\_CLOSED assemblies 1 norm survivors 1 final 0; Q 54 VERIFIED\_CLOSED assemblies 7 norm survivors 1 final 0; Q 63 VERIFIED\_CLOSED assemblies 15 norm survivors 1 final 0; Q 72 VERIFIED\_CLOSED assemblies 58 nor... \\
\texttt{verify\_mixed\_small} & mixed Q 60 VERIFIED\_WITH\_CITED\_CODE\_BOUND states 36 norm 3 residual 0; mixed Q 69 VERIFIED\_WITH\_CITED\_CODE\_BOUND states 1008 norm 26 residual 0; mixed Q 78 VERIFIED\_WITH\_CITED\_CODE\_BOUND states 17424 norm 201 residual 0 \\
\texttt{verify\_q81} & Q81 VERIFIED\_WITH\_CITED\_CODE\_BOUND isolate states 22 -> 0 non-isolate spectra 4 ->0 \\
\texttt{verify\_q99\_independent} & Q99 local rederivation PASS: 144,225,729 solutions; all residual survivors zero. \\
\texttt{verify\_component\_bounds} & Q 105 (14,1) no-isolate: 1 / 1 allocations; max/B 1.0; Q 114 (14,1) no-isolate: 2 / 2 allocations; max/B 0.9620720360256458; Q 123 (14,1) no-isolate: 3 / 3 allocations; max/B 0.9566417785303339; Q 132 (14,1) no-isolate: 4 / 4 allocations; max/B 0.9579734857152806; Q 141 (14,1) no-isolate: 5 / 5 allocations; max/B 0.9636743412482026; Q 150 (14,1) no-isolate: 6 / 6 allocations; max/B 0.9727713273053784; Q 159 (14,1) no-isolate: 7 / 7 allocations; max/B 0.9846300571620813; Q 168 (14,1) no-isolate: 8 / 8 allocations... \\
\texttt{audit\_all\_shells} & Global coverage: 19 closed; 19 unverified. Global maximality: False \\
\end{longtable}
\normalsize
All listed modules returned exit code 0 in the reconstructed replay. The aggregate shell audit reports 19 closed shells, 19 unverified shells, and \texttt{global\_maximality\_verified = false}. This is intentionally treated as an incomplete global proof.


The immutable-source manifest and the full delivered snapshot manifest have different purposes. A replay updates log timings and result files; the source manifest remains the pre-execution check. The original upload hash is in \file{source_provenance.json}. No complete repository checkout or missing remote proof file is implicitly included in these manifests.

\begin{thebibliography}{9}
\bibitem{TB} T. Todorov and G. Bogdanova. \emph{Ternary equidistant codes of length 11 through 15}. Journal of Mathematical and Computational Science 10 (2020), 2713--2721. Proposition 14 and Table 1. \url{https://doi.org/10.28919/jmcs/4964}.
\bibitem{repo} Order-15 research materials in \emph{order55-binary-circulant-maxdet}, owner \texttt{dongxuelian2}. Examined main ref \texttt{827672a2c60dad7064c7955d9c460261d866d25e} and research ref \texttt{007af7d8f7f644396aee27dd87c68c8f4169efd5}. Paths and recovery limits are documented in \file{recovered/source_notes.md}.
\bibitem{upload} User-supplied \file{order15_mu3_Q87-150_bundle.zip}. Sixteen recovered files, including the final ledger, proof digest, process notes and standalone $Q=99$ verifier. Preserved without adopting unsupported completion labels.
\end{thebibliography}
\end{document}
